\documentclass[11pt]{article}
\usepackage[margin=.8in,top=.7in,bottom=.7in]{geometry}
\usepackage{amsmath,amssymb,amsthm,hyperref}
\renewcommand{\normalsize}{\fontsize{10}{12}\selectfont}
\DeclareMathSizes{10}{10}{8}{6}
\hypersetup{colorlinks=true,linkcolor=blue,urlcolor=blue,citecolor=blue}
\setlength{\emergencystretch}{2em}
\newtheorem{theorem}{Theorem}
\newtheorem{lemma}{Lemma}
\newcommand{\A}{A}
\newcommand{\G}{G}
\begin{document}
\begin{center}
{\Large\bfseries The sharp top-mean index\\in dimensions at least five}\\[.5em]
Research proof and formalization packet\\
October 26, 2026
\end{center}
\begin{abstract}
For a probability vector in dimension $d$, form all $d^2$ arithmetic
means and all $d^2$ geometric means of its ordered coordinate pairs.
Sort the two arrays independently. We prove that, for every $d\ge5$,
the sum of their respective largest $k$ entries is at most $d$ for
every probability vector if and only if $k\le2d-5$.
This establishes the value proposed by Caro and Graswald in that
dimension range. The proof combines two incidence bounds for arbitrary
sets of ordered pairs with Cauchy--Schwarz. An explicit probability
vector proves sharpness in every dimension. The complete statement,
including the independently sorted arrays, is certified in Lean.
\end{abstract}

\section{Statement and source}
Let $p=(p_1,\ldots,p_d)$ satisfy $p_i\ge0$ and $\sum_i p_i=1$.
Set
\[
 a_{ij}=\frac{p_i+p_j}{2},\qquad g_{ij}=\sqrt{p_i p_j}
 \qquad(1\le i,j\le d).
\]
Write $a^\downarrow_1,\ldots,a^\downarrow_{d^2}$ and
$g^\downarrow_1,\ldots,g^\downarrow_{d^2}$ for these two arrays in
nonincreasing order, retaining repetitions, and put
\[
 A_k(p)=\sum_{\ell=1}^k a^\downarrow_\ell,
 \qquad G_k(p)=\sum_{\ell=1}^k g^\downarrow_\ell.
\]
The selected ordered pairs may differ between the two sums.
Caro--Graswald \cite{CG}, Conjecture~4.19 in arXiv~v3
(Conjecture~IV.19 in the published paper), propose the maximal index
$h(d)=2d-5$ for the universal inequality $A_{h(d)}+G_{h(d)}\le d$.
Their Remark~4.18 explains its role in necessary criteria for
infinitesimal Markovian divisibility. We prove the proposed value for
every $d\ge5$.

\begin{theorem}\label{main}
Let $d\ge5$ and $0\le k\le d^2$ be integers. Then
\[
 \bigl[A_k(p)+G_k(p)\le d\text{ for every probability vector }p\bigr]
 \quad\Longleftrightarrow\quad k\le2d-5.
\]
Moreover, $A_{2d-5}(p)+G_{2d-5}(p)<d$ for every such $p$.
\end{theorem}

For a nonnegative finite array, its largest-$k$ sum is the maximum of
the sums over $k$ distinct indices. Thus upper bounds may be proved
for arbitrary sets of ordered pairs, and lower bounds by exhibiting
such sets. This also handles ties and zero entries.

\section{The universal bound}
Put $m=d-3$ and $k=2m+1$. Write $P=\max_i p_i$.
For any set $S$ of $k$ ordered pairs define
\[
 c_i=\frac{\#\{j:(i,j)\in S\}+\#\{j:(j,i)\in S\}}2.
\]
Then the corresponding arithmetic sum is $\sum_i c_i p_i$, and
\begin{equation}\label{inc}
 c_i\le m+1,\qquad c_i+c_j\le m+\frac52\quad(i\ne j).
\end{equation}
Indeed, for one vertex the loop contributes $1$ and each other
incident ordered pair contributes $1/2$. For two distinct vertices,
at most four selected ordered pairs have both endpoints among them;
these contribute $1$ to $c_i+c_j$, and every other pair contributes
at most $1/2$. No symmetry assumption on $S$ is used.

Choose $u$ maximizing $c_i$, and let $a=c_u$ and
$b=\max_{i\ne u}c_i$. Since $a\ge b$, the arithmetic sum is at most
\[
 a p_u+b(1-p_u)\le aP+b(1-P).
\]
If $P\le2/3$, this is at most
\begin{equation}\label{lowA}
 \frac{2a+b}{3}=\frac{a+(a+b)}3\le\frac{4m+7}{6}.
\end{equation}
Also $\sum_{i,j}g_{ij}^2=(\sum_i p_i)^2=1$, so
Cauchy--Schwarz gives $G_k\le\sqrt{k}$. Consequently
\[
 A_k+G_k\le\frac{4m+7}{6}+\sqrt{2m+1}<m+3.
\]
For the strict inequality, the positive margin
$(2m+11)/6$ exceeds $\sqrt{2m+1}$ because
\[
 (2m+11)^2-36(2m+1)=(2m-7)^2+36>0.
\]

Now suppose $P>2/3$, and choose $u$ with $p_u=P$. Put
$\varepsilon=1-P<P/2$. Every arithmetic entry outside row and
column $u$ is at most $\varepsilon$. Give each selected entry the
baseline $P/2$. The loop $(u,u)$ needs an extra $P/2$; an entry
$(u,i)$ or $(i,u)$ with $i\ne u$ needs an extra $p_i/2$.
The loop appears at most once, and each tail coordinate appears in
at most two such entries. Hence
\begin{equation}\label{highA}
 A_k\le\frac{kP}{2}+\frac P2+\varepsilon=1+mP.
\end{equation}
The largest geometric entry is $g_{uu}=P$. A maximizing geometric
selection can include it: if necessary replace any selected index
by $(u,u)$, which does not decrease the sum. Applying
Cauchy--Schwarz to the remaining $k-1=2m$ entries gives
\begin{equation}\label{highG}
 G_k\le P+\sqrt{2m}\sqrt{1-P^2}.
\end{equation}
Combining \eqref{highA} and \eqref{highG}, and applying
two-dimensional Cauchy--Schwarz, yields
\begin{align*}
 A_k+G_k
 &\le1+(m+1)P+\sqrt{2m}\sqrt{1-P^2}\\
 &\le1+\sqrt{(m+1)^2+2m}
   =1+\sqrt{m^2+4m+1}<m+3.
\end{align*}
The final strict inequality follows from
$(m+2)^2-(m^2+4m+1)=3$. This proves the universal assertion at
$k=2d-5$. The same proof works for $d\ge3$.

\section{A uniform sharpness family}
Let $m=d-3\ge2$, $r=2m$, and $D=r^2+m$. Take
\[
 p=\frac1D\bigl(r^2,\underbrace{1,\ldots,1}_{m\text{ coordinates}},0,0\bigr).
\]
At the next index $k'=2m+2$, choose the following arithmetic entries:
the distinguished diagonal, all $2m$ orientations joining the
distinguished coordinate to the positive tail, and one orientation
joining it to a zero coordinate. Choose the geometric entries by
replacing this last arithmetic index with one positive-tail diagonal.
Each selection has exactly $2m+2$ distinct ordered pairs. Their sums are
\[
 A_{\rm sel}=\frac{(m+\tfrac32)r^2+m}{D},\qquad
 G_{\rm sel}=\frac{r^2+2mr+1}{D}.
\]
Since the independently sorted sums dominate these respective
selections,
\begin{align*}
 A_{k'}+G_{k'}-d
 &\ge\frac{-r^2/2+2mr-m^2-2m+1}{r^2+m}\\
 &=\frac{(m-1)^2}{4m^2+m}>0.
\end{align*}
For example, $d=5$ gives $p=(16,1,1,0,0)/18$ and selected sums
$29/9$ and $11/6$, whose total is $91/18>5$.
All entries are nonnegative, so both largest-$k$ sums are
nondecreasing in $k$. The witness excludes every $k\ge2d-4$;
the universal bound proves every $k\le2d-5$. This completes the
proof of Theorem~\ref{main}.

\section{Formal certificate and scope}
The certificate comprises 76 declarations in twelve Lean modules,
checked with Lean~4.34.1 and a pinned mathlib revision. Its endpoint
\texttt{MaxIndex.source\_index\_classification} is precisely the
equivalence in Theorem~\ref{main}, quantified over all
nonnegative functions $p:\mathrm{Fin}\,d\to\mathbb R$ of sum $1$.
The strict sufficiency theorem also covers $d\ge3$.
The development proves the literal sorted-array/finite-selection
equivalence, the loop and pair incidence counts, the geometric
square-mass identity, the largest-entry exchange, and the sizes and
sums of both sharpness selections. Reindexing the explicit witness
onto $\mathrm{Fin}\,d$ is proved as well. No combinatorial or
numerical certificate remains as an extra hypothesis.

The strict replay checks every declaration type and allows only
\texttt{propext}, \texttt{Classical.choice}, and \texttt{Quot.sound}.
The elementary incidence and Cauchy--Schwarz ingredients are
standard; the source-question conclusion is the combined sharp
index theorem for $d\ge5$. Small-dimensional classifications are
outside this certificate's maximality statement. Source and
bounded literature-search records accompany the development;
they do not certify absolute worldwide priority.

\begin{thebibliography}{9}
\bibitem{CG}
M.~C.~Caro and B.~R.~Graswald,
\emph{Necessary criteria for Markovian divisibility of linear maps},
J.~Math.~Phys. \textbf{62}, 042203 (2021),
\href{https://doi.org/10.1063/5.0031760}{doi:10.1063/5.0031760};
\href{https://arxiv.org/abs/2009.06666v3}{arXiv:2009.06666v3} (2022),
Conjecture~4.19.
\end{thebibliography}
\end{document}
